%% ============================================================
%%  Chapter 2 — Background: Halide Perovskite Semiconductors
%%              for Radiation Detection
%%  Dissertation: Radiation Responsive Materials for Advanced
%%  Radiation Detection and Systems
%% ============================================================

\chapter{Background: Halide Perovskite Semiconductors for Radiation Detection}

The research reported in Chapters~\ref{ch:CsFAPbI}--\ref{ch:MHyPbCl3}
rests on three intersecting bodies of knowledge: the crystal chemistry and
electronic structure of halide perovskites, the physics governing how
ionizing radiation deposits energy in matter, and the device physics that
determines how efficiently that deposited energy is converted into a
measurable electrical signal.  This chapter develops each body of knowledge
in sufficient depth to motivate the experimental choices made in later
chapters and to place the results in the context of the broader literature
on semiconductor radiation detectors.  The FPGA-based systems work is
treated separately in Chapter~\ref{ch:FPGA_background}.

%% ------------------------------------------------------------
\section{The Perovskite Crystal Structure and Material Families}
\label{sec:perovskite_structure}
%% ------------------------------------------------------------

\subsection{The ABX$_3$ Archetype}

The term ``perovskite'' derives from the mineral calcium titanate
(CaTiO$_3$), named for the Russian mineralogist Lev von Perovski and
first synthesized in the nineteenth century~\cite{Wells1945}.  The
defining structural motif is a three-dimensional network of corner-sharing
BX$_6$ octahedra in which a larger A-site cation occupies the
12-coordinate cuboctahedral void at the unit-cell body-center (in the
ideal cubic setting, space group $Pm\bar{3}m$).  The versatility of this
archetype arises from the wide range of chemistries that can be
accommodated without destroying the framework: in the oxide family,
B-sites range from Ti$^{4+}$ to Pb$^{4+}$, while in the halide family
addressed here, B is almost exclusively Pb$^{2+}$ (or Sn$^{2+}$) and
X is a monovalent halide (I$^{-}$, Br$^{-}$, or Cl$^{-}$).

The stability of a given ABX$_3$ composition is conveniently predicted by
the Goldschmidt tolerance factor,
\begin{equation}
  t = \frac{r_{\rm A} + r_{\rm X}}{\sqrt{2}\,(r_{\rm B} + r_{\rm X})},
  \label{eq:tolerance}
\end{equation}
where $r_i$ denotes the ionic radius of species $i$~\cite{Goldschmidt1926}.
Perovskite formation is generally favored for $0.80 \le t \le 1.05$;
values close to unity favor the ideal cubic structure, while smaller
values drive cooperative octahedral tilting that lowers the symmetry to
tetragonal or orthorhombic phases.  Kieslich et al.\ extended this
framework to hybrid organic-inorganic systems by defining effective radii
for polyatomic organic cations from the root-mean-square displacement of
their hydrogen envelope~\cite{Kieslich2015}, enabling quantitative phase
stability predictions for the methylammonium and formamidinium compounds
central to this work.

A complementary criterion is the octahedral factor $\mu = r_{\rm B}/r_{\rm X}$,
which must satisfy $0.414 \le \mu \le 0.732$ for the B cation to fit stably
inside the halide octahedron.  Pb$^{2+}$ ($r = 1.19$~\AA) paired with
I$^{-}$ ($r = 2.20$~\AA) gives $\mu = 0.54$, comfortably within the stable
range and consistent with the broad tunability of the Pb-halide framework.

\subsection{Three-Dimensional Structures}

\textbf{Methylammonium lead iodide (MAPbI$_3$).}
The most extensively studied hybrid halide perovskite, MAPbI$_3$
(methylammonium = CH$_3$NH$_3^+$, MA$^+$), adopts a cubic structure
above 327~K, a tetragonal structure between 162~K and 327~K, and an
orthorhombic structure below 162~K~\cite{Baikie2013}.  At room temperature,
the tetragonal phase (space group $I4/mcm$) is therefore the operative
structure for detector applications.  MAPbI$_3$ has a direct bandgap of
approximately 1.55~eV and a high optical absorption coefficient
($>10^4$~cm$^{-1}$ above the bandgap), along with the excellent charge
transport described in Section~\ref{sec:hecht}.  Large single crystals
with centimeter-scale dimensions can be grown from solution by inverse
temperature crystallization~\cite{Saidaminov2015}, making bulk detector
fabrication feasible.  The principal weakness of MAPbI$_3$ for detector
applications is the light A-site cation: the low atomic masses of C, N,
and H contribute minimally to gamma-ray stopping power, and the modest
hydrogen density in the cubic CH$_3$NH$_3^+$ cation limits fast-neutron
response.  The MAPbI$_3$-hBN composite detector described in
Chapter~\ref{ch:MAPbI3hBN} addresses the neutron sensitivity limitation
by incorporating high-\emph{n} boron-containing nanoflakes within a
MAPbI$_3$ matrix.

\textbf{Cesium lead bromide (CsPbBr$_3$).}
Replacing the organic MA$^+$ cation with the monoatomic Cs$^+$
($r = 1.67$~\AA) yields an all-inorganic perovskite with significantly
improved thermal and photochemical stability~\cite{Stoumpos2013}.
CsPbBr$_3$ undergoes orthorhombic-to-tetragonal ($\sim$88~$^\circ$C)
and tetragonal-to-cubic ($\sim$130~$^\circ$C) transitions, remaining
in the orthorhombic phase at room temperature.  Its bandgap of
approximately 2.3~eV is wide enough to suppress intrinsic carrier
excitation at 300~K, yielding resistivities exceeding $10^9$~$\Omega$\,cm
in single crystals~\cite{He_NatPhoton_2021} — a prerequisite for the low
dark-current operation required in radiation counting.  Most critically
for gamma detection, the presence of both Cs ($Z = 55$) and Pb ($Z = 82$)
gives CsPbBr$_3$ a calculated effective atomic number of approximately
65.9, far exceeding that of conventional detector materials such as silicon
($Z_{\rm eff} = 14$) and approaching that of cadmium zinc telluride
($Z_{\rm eff} \approx 49.1$)~\cite{Springer_CsPbBr3}.  This high
$Z_{\rm eff}$ makes CsPbBr$_3$ among the most promising semiconductor
candidates for room-temperature gamma spectroscopy.

\textbf{Mixed cation Cs$_{0.1}$FA$_{0.9}$PbI$_3$.}
Pure formamidinium lead iodide (FAPbI$_3$, FA$^+$ = CH(NH$_2$)$_2^+$)
has a narrower bandgap of approximately 1.48~eV and superior thermal
stability relative to MAPbI$_3$, but the desired black
$\alpha$-phase is metastable at room temperature, spontaneously converting
to the photoinactive yellow $\delta$-phase~\cite{Weller2015}.  Partial
substitution of FA$^+$ with Cs$^+$ (typically 5--15 mol\%) suppresses
this phase transition by reducing the tolerance factor toward unity,
entropy-stabilizing the cubic $\alpha$-phase at ambient
conditions~\cite{Yi2016}.  The composition Cs$_{0.1}$FA$_{0.9}$PbI$_3$
investigated in Chapter~\ref{ch:CsFAPbI} exploits this stabilization
while retaining the narrow bandgap and large absorption cross-section of
FAPbI$_3$, enabling its use as a photon-responsive active layer in a
photovoltaic architecture repurposed as a radiation monitor.

\subsection{Two-Dimensional Ruddlesden-Popper Structures}

The three-dimensional perovskite framework can be sliced along specific
crystallographic planes and the resulting slabs separated by layers of
large organic spacer cations to form two-dimensional (2D) or quasi-2D
structures.  In the Ruddlesden-Popper series the general formula is
$(R)_2$(MA)$_{n-1}$Pb$_n$X$_{3n+1}$, where $R$ is a large organic
cation that cannot fit the A-site of the 3D structure, $n$ is the number
of corner-sharing PbX$_6$ octahedral layers between adjacent organic
spacer bilayers, and the ratio of inorganic to organic content increases
with $n$~\cite{Tsai2018}.  In the limit $n \rightarrow \infty$ the
three-dimensional structure is recovered.

For the present work, the relevant composition is the $n = 3$ member
PEA$_2$MA$_2$Pb$_3$I$_{10}$, where PEA$^+$ denotes the
phenylethylammonium cation (C$_6$H$_5$(CH$_2$)$_2$NH$_3^+$).
The alternating organic-inorganic layering creates a natural quantum-well
electrostatic landscape in which charge is preferentially confined within
the higher-dielectric inorganic slabs.  From a detector perspective, this
layered architecture confers two important advantages over its 3D
counterpart: the bulky PEA$^+$ layers act as physical barriers to halide
vacancy migration, substantially suppressing the ion-transport instabilities
detailed in Section~\ref{sec:ion_migration}; and the organic layers provide
intrinsic passivation of the inorganic slab surfaces, reducing
trap-state densities at the heterojunction interfaces~\cite{Cao2018}.
The price paid is a modest reduction in charge-carrier mobility relative
to three-dimensional analogs, since carriers must traverse the
organic-inorganic boundaries.  Chapter~\ref{ch:2Dperovskite} investigates
whether surface-passivation strategies can recover sufficient transport
quality to enable practical detector operation while retaining the
stability advantages of the 2D architecture.

\subsection{Methylhydrazinium Lead Chloride (MHyPbCl$_3$)}

Methylhydrazinium lead chloride, MHyPbCl$_3$ (MHy$^+$ =
CH$_3$NH$_2$NH$_2^+$), is a relatively recently reported member of the
halide perovskite family in which the conventional monovalent cation is
replaced by the methylhydrazinium ion~\cite{Nowok2021}.  This composition
is of particular interest for fast-neutron detection because of its
exceptionally high hydrogen content: the CH$_3$NH$_2$NH$_2^+$ cation
carries five hydrogen atoms, and the resulting crystallographic hydrogen
density is approximately $3.79 \times 10^{22}$~H\,cm$^{-3}$~\cite{Panaccione_NIM_2024}.
This exceeds the hydrogen density of many purpose-built neutron moderator
and detector materials.  Fast neutrons (energies above $\sim$1~MeV)
interact primarily with hydrogen nuclei via elastic scattering, transferring
up to 100\% of their kinetic energy to the recoil proton in a head-on
collision; those recoil protons then ionize the detector medium in the
same manner as a directly ionizing particle.  The combination of high
hydrogen density (for neutron moderation and proton production), moderate
bandgap ($\sim$3.0~eV, appropriate for room-temperature operation), and
the semiconducting perovskite framework (for direct electronic readout)
makes MHyPbCl$_3$ a compelling candidate for a compact direct-conversion
fast-neutron detector.  Chapter~\ref{ch:MHyPbCl3} presents the synthesis,
characterization, and fast-neutron response measurements for this material.

%% ------------------------------------------------------------
\section{Radiation Detection Fundamentals}
\label{sec:detection_fundamentals}
%% ------------------------------------------------------------

\subsection{Mechanisms of Ionizing Radiation Detection}

A radiation detector converts the energy deposited by an ionizing particle
or photon into a measurable quantity — in semiconductor detectors, this
quantity is an electrical charge pulse whose amplitude is proportional to
the deposited energy.  The process proceeds in three stages: (1) the
primary ionizing event, in which the incident radiation transfers energy
to the detector medium, producing a spatial distribution of electron-hole
pairs along the particle track or at the absorption site; (2) charge
transport, in which the liberated carriers drift under the applied bias
field toward the collection electrodes; and (3) signal induction and
readout, in which the moving charges induce a current pulse on the
electrodes that is amplified and digitized~\cite{knoll2010}.

The mean energy required to create one electron-hole pair in a
semiconductor is given by Klein's empirical rule,
\begin{equation}
  W \approx 2.67\,E_{\rm g} + 0.87~\text{eV},
  \label{eq:klein}
\end{equation}
where $E_{\rm g}$ is the bandgap energy~\cite{Klein1968}.  For
MAPbI$_3$ ($E_{\rm g} \approx 1.55$~eV), this gives $W \approx 5.0$~eV
per pair; for CsPbBr$_3$ ($E_{\rm g} \approx 2.3$~eV), $W \approx 7.0$~eV
per pair.  These values compare favorably with conventional semiconductor
detectors: silicon ($W = 3.6$~eV) and germanium ($W = 2.9$~eV) produce
more pairs per unit energy, but those materials require cryogenic cooling
or high-purity growth technologies incompatible with cost-effective
deployment.  The higher $W$ of perovskites is partially offset by their
far greater stopping power per unit thickness (Section~\ref{sec:radiation_interaction}),
so the total charge produced by a given photon energy can be comparable
to or exceed that in a thinner silicon detector despite the higher pair
creation energy.

\subsection{Energy Resolution and Pulse Height Analysis}

The ability to resolve closely spaced photon energies — and thereby identify
specific radionuclides — is quantified by the energy resolution $R$,
conventionally expressed as the fractional full-width at half-maximum (FWHM)
of a spectral peak:
\begin{equation}
  R = \frac{\Delta E_{\rm FWHM}}{E_0},
  \label{eq:resolution}
\end{equation}
where $E_0$ is the incident photon energy.  In an ideal semiconductor
detector in which Poisson statistics govern pair production, the minimum
achievable resolution is set by the Fano limit:
\begin{equation}
  \Delta E_{\rm FWHM}^{\rm Fano} = 2.35\,\sqrt{F \cdot W \cdot E_0},
  \label{eq:fano}
\end{equation}
where $F$ is the Fano factor ($F \approx 0.1$--0.2 for most semiconductors)~\cite{knoll2010}.
In practice, the measured resolution is degraded by incomplete charge
collection due to carrier trapping, electronic noise from the amplifier chain,
and dielectric polarization from ion migration (the last factor being
specific to halide perovskites and discussed in Section~\ref{sec:ion_migration}).
The benchmark figure commonly cited in the perovskite detector literature is
the 1.4\% resolution at 662~keV ($^{137}$Cs) achieved by He et al.\ for
a CsPbBr$_3$ single crystal coupled to a custom application-specific
integrated circuit (ASIC)~\cite{He_NatPhoton_2021}, which at the time
of publication was the best room-temperature resolution demonstrated by
any non-CZT semiconductor detector.

Energy-resolved radiation measurements are performed by connecting the
detector preamplifier output to a pulse-height analyzer (PHA), which bins
each charge pulse by its amplitude and accumulates a histogram — the
pulse-height spectrum — whose peak positions encode incident photon
energies and whose areas are proportional to the detected activity.
The shaping time of the amplifier is a critical parameter: it must be
long enough to integrate the slow hole component of the charge pulse
(hole mobilities in halide perovskites are typically 5--10$\times$ lower
than electron mobilities) but short enough to avoid excessive pileup at
practical count rates.

\subsection{X-Ray Detector Sensitivity}

For applications in medical and industrial X-ray imaging, the relevant
figure of merit is not energy resolution but sensitivity $S$, defined as
the charge collected per unit area per unit X-ray exposure:
\begin{equation}
  S = \frac{\Delta Q}{A \cdot \Delta X} \quad \left[\mu\text{C}\,\text{Gy}^{-1}\,\text{cm}^{-2}\right],
  \label{eq:sensitivity}
\end{equation}
where $\Delta X$ is the absorbed dose rate.  Sensitivity is controlled by
the X-ray absorption efficiency (determined by the stopping power and
detector thickness), the charge collection efficiency (determined by $\mu\tau$
products and applied field), and the dark current level (which sets the
minimum detectable signal).  Perovskite thick-film detectors fabricated
from polycrystalline MAPbI$_3$ or MAPbBr$_3$ have demonstrated sensitivities
exceeding 4,500~$\mu$C\,Gy$^{-1}$\,cm$^{-2}$ under 8~keV X-rays,
more than $10^4$ times higher than conventional amorphous selenium
panels~\cite{Wei2017}.  The high atomic number of the Pb sublattice is
the primary driver of this extraordinary sensitivity at soft X-ray
energies, where the photoelectric cross-section scales roughly as $Z^{4.5}$.

\subsection{Neutron Detection and Discrimination}

Neutron detection in semiconductor materials is indirect: neutrons carry
no charge and interact negligibly with electrons, so they must first
undergo a nuclear reaction or elastic scattering event that produces a
charged secondary particle capable of ionizing the detector.  Thermal
neutrons (energies below $\sim$0.5~eV) are most efficiently captured via
exothermic reactions such as $^{10}$B(n,$\alpha$)$^7$Li and
$^6$Li(n,$\alpha$)$^3$H.  Fast neutrons above $\sim$1~MeV interact
principally through elastic scattering from hydrogen nuclei, producing
recoil protons with a continuous energy spectrum up to the full neutron
energy.  The detection efficiency for fast neutrons is proportional to
the hydrogen areal density in the active medium, making hydrogen-rich
materials intrinsically advantageous.  The measured $^{1}$H density in
MHyPbCl$_3$ of $3.79 \times 10^{22}$~cm$^{-3}$ exceeds that of common
polymer moderator materials and places this perovskite in a unique
position as a combined moderating and detecting medium for fast
neutrons~\cite{Panaccione_NIM_2024}.

Discrimination between neutron and gamma signals is accomplished
by exploiting differences in the ionization density (linear energy
transfer, LET) of the respective secondary particles.  Recoil protons
from fast-neutron elastic scattering and alpha particles from thermal
capture reactions deposit energy far more densely than the Compton
electrons produced by gamma photons of similar incident energy.  In
systems with sufficient time resolution, pulse-shape discrimination
(PSD) can separate the fast and slow scintillation components that
arise from these different ionization densities; in solid-state detectors,
differences in carrier extraction dynamics under the high local ionization
density of heavy charged particles can be exploited in a similar
manner~\cite{knoll2010}.

%% ------------------------------------------------------------
\section{Wide Bandgap Semiconductors for Room-Temperature Detection}
\label{sec:wide_bandgap}
%% ------------------------------------------------------------

\subsection{The Room-Temperature Constraint}

The operating requirement most often overlooked in academic detector
comparisons — yet most consequential for deployment in nuclear plants,
field instruments, and space systems — is that the detector must
function reliably at room temperature without active cooling.  This
single constraint eliminates two of the highest-performing semiconductor
detector technologies: germanium (HPGe) must be operated below 120~K
to suppress thermally generated leakage current, requiring liquid-nitrogen
dewars or mechanical cryocoolers that add cost, mass, and failure modes
incompatible with unattended monitoring~\cite{knoll2010}.  Silicon drift
detectors, while operable at modest Peltier cooling ($-20$ to
$-40~^\circ$C), still suffer from displacement damage under high neutron
fluences that degrade their spectroscopic performance irreversibly.

The root cause of the temperature sensitivity in narrow-gap semiconductors
is the intrinsic carrier concentration $n_i$, which increases
exponentially with temperature as:
\begin{equation}
  n_i = \sqrt{N_c N_v}\,\exp\!\left(-\frac{E_{\rm g}}{2kT}\right),
  \label{eq:ni}
\end{equation}
where $N_c$ and $N_v$ are the effective density of states in the
conduction and valence bands, respectively, and $k$ is Boltzmann's
constant~\cite{Sze2007}.  At room temperature ($T = 300$~K), germanium
($E_{\rm g} = 0.67$~eV) has $n_i \approx 2.4 \times 10^{13}$~cm$^{-3}$,
producing a leakage current density many orders of magnitude larger than
the signal currents from typical radiation fields.  Silicon
($E_{\rm g} = 1.12$~eV) fares considerably better with
$n_i \approx 1.5 \times 10^{10}$~cm$^{-3}$, but displacement damage
from neutron irradiation introduces compensating defects that effectively
raise the intrinsic carrier density and degrade the depletion voltage
over time.

For a semiconductor to operate as a room-temperature detector without
active cooling or impractically high reverse bias, $n_i$ must be
sufficiently low that the thermal generation current is smaller than the
smallest radiation-induced signal of interest.  A widely cited rule of
thumb requires $n_i \lesssim 10^7$~cm$^{-3}$, which for silicon-like
band structure translates to a minimum bandgap of approximately
1.5~eV~\cite{Owens2004}.  This threshold defines the wide bandgap
semiconductor regime relevant to room-temperature radiation detection.

\subsection{Figures of Merit for Wide Bandgap Detector Materials}

The suitability of a wide bandgap semiconductor for radiation detection
is characterized by several interrelated figures of merit beyond the
bandgap itself:

\textbf{Resistivity} ($\rho$, $\Omega$\,cm).  High bulk resistivity
minimizes the dark current $I_{\rm dark} = V/R = VA/(L\rho)$ that flows
under the applied bias field and sets the noise floor.  For
spectroscopic applications, resistivities above $10^9$~$\Omega$\,cm
are typically required; for imaging applications where high flux
dominates shot noise, values above $10^7$~$\Omega$\,cm may be
acceptable.  Wide bandgap materials achieve high resistivity by a
combination of large $E_{\rm g}$ (reducing $n_i$) and the absence of
shallow donor or acceptor impurities in the bandgap.  CsPbBr$_3$
single crystals have been measured at $\rho > 10^{10}$~$\Omega$\,cm
at room temperature~\cite{He_NatPhoton_2021}, exceeding the resistivity
of commercial CdZnTe by approximately one order of magnitude.

\textbf{Mobility-lifetime product} ($\mu\tau$, cm$^2$\,V$^{-1}$).
As derived in Section~\ref{sec:hecht}, $\mu\tau$ directly governs
the mean free path of carriers between generation and trapping events.
Wide bandgap semiconductors generally have lower carrier mobilities
than narrow-gap materials (because heavier effective masses accompany
larger bandgaps in many crystal systems), but the lifetime can be
extraordinarily long in materials with low defect densities, partially
compensating.  The target for spectroscopic detectors is
$\mu_e\tau_e \gtrsim 10^{-3}$~cm$^2$\,V$^{-1}$; for hole collection,
the equivalent target is often relaxed to $10^{-4}$~cm$^2$\,V$^{-1}$
when single-polarity collection schemes are employed~\cite{Luke1995}.

\textbf{Effective atomic number} ($Z_{\rm eff}$).  For gamma and X-ray
detection, $Z_{\rm eff}$ determines the photoelectric stopping power,
which scales as $Z^{4.5}$.  Wide bandgap semiconductors based on
heavy elements (Pb, Bi, Hg, Cd, Te) combine the thermal benefits of
a large bandgap with high stopping power in a way that narrow-gap
elemental semiconductors cannot achieve.  Silicon ($Z = 14$) requires
centimeters of active thickness to stop 662~keV photons, while a few
millimeters of CsPbBr$_3$ ($Z_{\rm eff} = 65.9$) or CdZnTe
($Z_{\rm eff} = 49.1$) provide equivalent absorption~\cite{Murty1965}.

\textbf{Pair creation energy} ($W$, eV).  As noted in
Eq.~\ref{eq:klein}, $W$ is approximately proportional to $E_{\rm g}$.
Wide bandgap materials therefore produce fewer electron-hole pairs per
unit of deposited energy than narrow-gap ones, increasing the
statistical (Fano) contribution to the energy resolution.  This is the
principal penalty paid for room-temperature operability: it must be
offset by large $\mu\tau$ products that ensure complete charge
collection, since any additional variance from trapping further
degrades the resolution beyond the Fano limit.

\subsection{The Wide Bandgap Material Landscape}

Table~\ref{tab:wbg_comparison} contextualizes the halide perovskites
investigated in this dissertation within the wider landscape of wide
bandgap semiconductor detector materials.

\begin{table}[htbp]
\centering
\caption{Selected properties of wide bandgap semiconductor detector
         materials at 300~K.  CZT = cadmium zinc telluride
         Cd$_{0.9}$Zn$_{0.1}$Te. Data compiled
         from~\cite{Owens2004,He_NatPhoton_2021,Springer_CsPbBr3,knoll2010}.}
\label{tab:wbg_comparison}
\startsinglespace
\begin{tabular}{lcccccc}
\hline
Material & $E_{\rm g}$ (eV) & $Z_{\rm eff}$ &
  $\rho$ ($\Omega$\,cm) & $\mu_e\tau_e$ (cm$^2$\,V$^{-1}$) &
  $W$ (eV) & Best $R_{662}$ \\
\hline
Si          & 1.12 & 14   & $>10^4$   & $>1$                  & 3.6  & $<0.1$\% (cryo) \\
CdTe        & 1.50 & 49.5 & $10^9$    & $3\times10^{-3}$      & 4.4  & 1.7\% \\
CZT         & 1.57 & 49.1 & $10^{10}$ & $10^{-2}$             & 4.6  & 0.5\% \\
GaAs        & 1.42 & 31.6 & $10^7$    & $10^{-5}$             & 4.2  & $\sim$2\% \\
4H-SiC      & 3.26 & 14   & $>10^{12}$& $10^{-6}$             & 8.0  & $\sim$1\%* \\
HgI$_2$     & 2.13 & 62.2 & $10^{13}$ & $10^{-4}$             & 4.2  & 0.5\% \\
CsPbBr$_3$  & 2.30 & 65.9 & $>10^{10}$& $1.3\times10^{-2}$   & 7.0  & 1.4\% \\
MAPbI$_3$   & 1.55 & 51.4 & $10^8$    & $8\times10^{-3}$      & 5.0  & $\sim$6\% \\
\hline
\multicolumn{7}{l}{\footnotesize *For alpha particles; gamma spectroscopy
  performance is limited by the low $Z_{\rm eff}$.}
\end{tabular}
\end{table}

Several observations emerge from this comparison.  CdZnTe has long
occupied the upper tier of room-temperature gamma spectrometers, combining
a near-ideal bandgap for thermal leakage suppression, high $Z_{\rm eff}$,
and a $\mu_e\tau_e$ product sufficient for millimeter-to-centimeter-scale
crystals — but its dominance rests entirely on the extraordinary effort
invested in Bridgman crystal growth and post-growth processing that has
accumulated over three decades.  HgI$_2$ offers the highest $Z_{\rm eff}$
and resistivity of the conventional materials but suffers from mechanical
fragility and a phase transition at 127~$^\circ$C that has limited
commercial uptake.  Silicon carbide (4H-SiC) is exceptional in its
radiation hardness and thermal tolerance but its low $Z_{\rm eff}$
confines it to charged-particle and neutron applications where stopping
power from the heavy-element sublattice is not required~\cite{Owens2004}.

CsPbBr$_3$ occupies a remarkably favorable position in this landscape.
Its $Z_{\rm eff}$ of 65.9 is the highest of any material in common
laboratory study, its $\mu_e\tau_e$ is competitive with CdZnTe despite
far simpler solution-based crystal growth, and its resistivity exceeds
that of CdZnTe by an order of magnitude.  The 1.4\% energy resolution
at 662~keV achieved by He et al.\ is already within a factor of three
of the best CdZnTe single-element detectors, and it was achieved after
only a few years of focused development~\cite{He_NatPhoton_2021}.  The
remaining gap is attributable primarily to the ion-migration polarization
effect (Section~\ref{sec:ion_migration}), which introduces a
time-dependent component to charge collection efficiency not present in
purely electronic semiconductors.

\subsection{Why Bandgap Engineering Matters: The Perovskite Advantage}

The halide perovskite family is unique among wide bandgap detector
candidates in that the bandgap is continuously tunable over the range
1.2--3.5~eV by compositional substitution at all three crystallographic
sites, without changing the crystal structure or growth method.  Replacing
iodide with bromide shifts the bandgap from 1.55~eV (MAPbI$_3$) to
2.3~eV (MAPbBr$_3$); substituting Pb$^{2+}$ with Sn$^{2+}$ narrows the
gap to $\sim$1.2~eV; replacing MA$^+$ with the larger FA$^+$ cation
narrows the gap marginally to $\sim$1.48~eV~\cite{Stoumpos2013}.
This tunability enables the optimization of three competing
requirements simultaneously: (1) $E_{\rm g}$ large enough to suppress
thermal leakage at the maximum operating temperature; (2) $E_{\rm g}$
small enough to keep $W$ — and therefore the Fano noise — acceptably
low; and (3) $Z_{\rm eff}$ high enough for efficient photon stopping
in the energy range of interest.  No other material family offers all
three degrees of freedom within a single, solution-processable crystal
system, which is the central motivation for investigating multiple
perovskite compositions across the chapters that follow.

%% ------------------------------------------------------------
\section{Radiation Interaction with High-\textit{Z} Semiconductors}
\label{sec:radiation_interaction}
%% ------------------------------------------------------------

\subsection{Photon Interactions}

Three mechanisms govern the interaction of gamma and X-ray photons with
matter: the photoelectric effect, Compton scattering, and pair production.
Each dominates over a different energy range and scales differently with
the atomic number of the absorber, making the choice of detector material
central to the detection efficiency for a given source energy.

The \emph{photoelectric effect} is dominant below approximately 100~keV
in high-$Z$ materials.  An incident photon is completely absorbed by an
inner-shell electron, which is ejected with kinetic energy $E_e = h\nu - E_b$,
where $E_b$ is the binding energy of that shell.  The photoelectric
cross-section scales strongly with atomic number and inversely with photon
energy as $\sigma_{\rm pe} \propto Z^{4.5}/E^{3.5}$, so high-$Z$
semiconductors provide orders-of-magnitude higher photoelectric absorption
than silicon or germanium for the same detector volume.  The complete
energy deposition of the photoelectric process produces the full-energy
peak in the pulse-height spectrum, making it the principal mechanism for
photon spectroscopy.

\emph{Compton scattering} dominates for photon energies between
approximately 100~keV and several MeV in heavy materials.  An incident
photon transfers a fraction of its energy to a loosely bound or free
electron, with the scattered photon carrying the remainder.  Because
only partial energy is deposited, Compton events populate a continuum
below the full-energy peak rather than contributing to it.  The Compton
cross-section is proportional to $Z$, so high-$Z$ absorbers still offer
a substantial geometric advantage, but the spectroscopic information is
contained only in the photopeak fraction.

\emph{Pair production} becomes significant above the threshold energy of
1.022~MeV and dominates at multi-MeV photon energies.  The photon is
converted into an electron-positron pair in the Coulomb field of the
nucleus; the 511-keV annihilation photons from the subsequent positron
annihilation appear as characteristic escape peaks in the spectrum.
The pair-production cross-section scales as $Z^2$.

The combined effect of these mechanisms is captured by the mass
attenuation coefficient $\mu/\rho$, tabulated as a function of energy
for all elements by the National Institute of Standards and Technology
(NIST XCOM database).  For a compound material, the effective atomic
number is calculated as:
\begin{equation}
  Z_{\rm eff} = \left(\sum_i w_i Z_i^n\right)^{1/n},
  \label{eq:Zeff}
\end{equation}
where $w_i$ is the weight fraction of element $i$ and the exponent
$n \approx 3.5$ is chosen to match the photoelectric cross-section
scaling~\cite{Murty1965}.  Applying this formula to CsPbBr$_3$
(Cs: $Z = 55$, Pb: $Z = 82$, Br: $Z = 35$; stoichiometric weight fractions
of 0.229, 0.449, and 0.322, respectively) yields $Z_{\rm eff} \approx 65.9$,
substantially higher than silicon ($Z_{\rm eff} = 14$), germanium (32),
or even CdZnTe ($Z_{\rm eff} \approx 49.1$)~\cite{Springer_CsPbBr3}.

\subsection{Charge Transport and the Hecht Equation}
\label{sec:hecht}

Even if radiation deposits charge in the bulk of a detector, that charge
must be transported to the collection electrodes without recombination
or trapping if a proportional signal is to be measured.  The quality of
charge transport is characterized by the mobility-lifetime product $\mu\tau$
for electrons and holes separately.  Carrier \emph{mobility} $\mu$
(cm$^2$\,V$^{-1}$\,s$^{-1}$) quantifies the drift velocity per unit
electric field; carrier \emph{lifetime} $\tau$ (s) is the mean time before
a carrier is irreversibly trapped at a deep defect and removed from the
circuit.  The product $\mu\tau$ has units of cm$^2$\,V$^{-1}$ and
represents the mean distance a carrier drifts per unit of applied field
before being trapped; a large $\mu\tau$ indicates that carriers can
traverse a thick crystal without significant charge loss.

The fractional charge collection for a detector of thickness $L$ under
applied field $E$ can be derived rigorously for a planar single-carrier
case.  For a general double-carrier detector with the interaction site
at depth $x$ from the cathode, the Hecht equation gives:
\begin{equation}
  Q(x) = Q_0 \left[
    \frac{\mu_e \tau_e E}{L}\left(1 - e^{-(L-x)/\mu_e\tau_e E}\right)
    +
    \frac{\mu_h \tau_h E}{L}\left(1 - e^{-x/\mu_h\tau_h E}\right)
  \right],
  \label{eq:hecht}
\end{equation}
where $Q_0 = eN_0$ is the charge generated by the initial ionizing
event~\cite{Hecht1932}.  For uniform irradiation, averaging
Eq.~\ref{eq:hecht} over $x$ from 0 to $L$ gives the mean charge
collection efficiency (CCE), which approaches unity only when both
$\mu_e\tau_e E/L \gg 1$ and $\mu_h\tau_h E/L \gg 1$.  In practice
this requires either very high applied fields, very thin detectors,
or very large $\mu\tau$ products.

Halide perovskites have demonstrated $\mu\tau$ products competitive with
those of mature compound semiconductors.  Single-crystal MAPbI$_3$
exhibits $\mu_e\tau_e \approx 8 \times 10^{-3}$~cm$^2$\,V$^{-1}$ and
$\mu_h\tau_h \approx 3 \times 10^{-3}$~cm$^2$\,V$^{-1}$~\cite{Saidaminov2015},
values that in a 1-mm-thick crystal at a modest field of 100~V\,cm$^{-1}$
yield a mean electron drift length of $\sim$8~mm — far exceeding the
crystal thickness.  CsPbBr$_3$ single crystals have achieved
$\mu_e\tau_e$ up to $1.34 \times 10^{-2}$~cm$^2$\,V$^{-1}$, supporting
the 1.4\% energy resolution reported by He et al.~\cite{He_NatPhoton_2021}.
These values are within a factor of two to three of CdZnTe
($\mu_e\tau_e \sim 10^{-2}$~cm$^2$\,V$^{-1}$), a material that has
required decades of crystal growth refinement to achieve its current
performance levels.

The hole transport is typically the limiting factor in halide perovskites,
as in most compound semiconductors.  The hole $\mu\tau$ in both MAPbI$_3$
and CsPbBr$_3$ is typically 5--10$\times$ lower than the electron value,
leading to a characteristic single-polarity collection signature in which
events occurring near the cathode (from which electrons must traverse the
full crystal thickness) are preferentially detected.  Electronic techniques
such as single-carrier collection (using a Frisch collar or virtual Frisch
ring) and waveform analysis can mitigate poor hole collection~\cite{Luke1995},
and perovskite detectors are amenable to these approaches.

%% ------------------------------------------------------------
\section{Schottky Diode Contacts in Semiconductor Detectors}
\label{sec:schottky}
%% ------------------------------------------------------------

\subsection{Metal-Semiconductor Junction Physics}

When a metal is brought into contact with a semiconductor, charge
redistributes at the interface until the Fermi levels align, creating
a built-in potential difference and a depletion region in the
semiconductor.  For a metal in contact with an n-type semiconductor,
the Schottky barrier height is:
\begin{equation}
  \phi_{\rm B} = \phi_m - \chi_s,
  \label{eq:schottky_barrier}
\end{equation}
where $\phi_m$ is the metal work function and $\chi_s$ is the electron
affinity of the semiconductor~\cite{Sze2007}.  The depletion region
extends a distance $W$ into the semiconductor:
\begin{equation}
  W = \sqrt{\frac{2\varepsilon\varepsilon_0 (V_{\rm bi} - V)}{e N_D}},
  \label{eq:depletion_width}
\end{equation}
where $\varepsilon$ is the static dielectric constant, $V_{\rm bi}$
is the built-in voltage, $V$ is the applied reverse bias, and $N_D$
is the donor concentration.  The rectifying current-voltage relationship
of the Schottky diode follows the thermionic emission model:
\begin{equation}
  J = J_0\left(e^{qV/nkT} - 1\right), \quad
  J_0 = A^{**}T^2\,e^{-q\phi_{\rm B}/kT},
  \label{eq:IV_schottky}
\end{equation}
where $A^{**}$ is the effective Richardson constant and $n$ is the
ideality factor~\cite{Sze2007}.

In contrast, an \emph{ohmic} contact is formed when the metal work
function is matched to the Fermi level of the semiconductor such that
no barrier to majority-carrier flow exists.  Ohmic contacts allow
free exchange of majority carriers with the external circuit and are
necessary on the collection side of the detector to avoid field
distortion near the electrode.

\subsection{Application to Halide Perovskite Detectors}

The asymmetric contact architecture — one ohmic and one Schottky
electrode — has become the standard configuration for halide perovskite
radiation detectors because it simultaneously achieves low dark current
and efficient charge collection~\cite{Wei2016}.  Under reverse bias, the
Schottky junction is depleted, suppressing majority-carrier injection
from the blocking contact; minority carriers generated by ionizing
radiation within the depletion region are swept by the electric field
toward the appropriate electrode.  Because the dark current is
thermally limited by the Schottky barrier rather than set by the bulk
resistivity, even moderate barrier heights of 0.5--0.8~eV can reduce
the dark current density by several orders of magnitude relative to the
ohmic-contact geometry.

For MAPbI$_3$-based detectors, gold (Au, $\phi_m = 5.1$~eV) forms a
near-ohmic contact on the p-type surface, while chromium (Cr,
$\phi_m = 4.5$~eV), gallium, or indium contacts on the opposite face
exhibit Schottky character~\cite{Wei2016}.  For CsPbBr$_3$, the typical
asymmetric pair is a carbon paste or Au ohmic contact on one face and
a Ga or indium-gallium eutectic Schottky contact on the other.  The
net result is a device structure functionally analogous to a $p$-$i$-$n$
junction in which the intrinsic region (the semi-insulating bulk crystal)
contains the active detector volume, and the contact asymmetry enforces
unipolar drift under bias~\cite{He_NatPhoton_2021}.

For the 2D perovskite devices studied in Chapter~\ref{ch:2Dperovskite},
the layered crystal architecture introduces an additional consideration:
charge transport perpendicular to the organic spacer layers (the
out-of-plane direction, which is the relevant direction for planar
detector geometry) is more resistive than in-plane transport, and the
contact interface with the topmost organic spacer layer must be carefully
engineered to avoid a secondary barrier that would impede collection.
Surface passivation using Lewis-base molecules that coordinate to
surface Pb$^{2+}$ sites can simultaneously reduce trap-state density
and tune the contact work function alignment.

%% ------------------------------------------------------------
\section{Ion Migration in Halide Perovskites}
\label{sec:ion_migration}
%% ------------------------------------------------------------

\subsection{Origins and Mechanisms}

The same structural softness that enables solution-processable
crystal growth from precursor salts also makes halide perovskites
susceptible to ionic transport under applied fields — a property that
distinguishes them fundamentally from conventional covalent or ionic
semiconductors such as silicon or CdZnTe.  Density functional theory
calculations and activation energy measurements converge on halide
vacancy migration as the dominant ionic conduction mechanism, with
activation energies of approximately 0.1--0.6~eV depending on the
halide and A-site cation~\cite{Eames2015}.  For comparison, the
thermal energy at room temperature ($kT \approx 26$~meV) is
sufficient to drive significant vacancy hopping across these barriers
on detector-relevant timescales, particularly under the electric fields
of 10--1,000~V\,cm$^{-1}$ routinely applied to semiconductor detectors.

The migration pathway for iodide vacancies in MAPbI$_3$ proceeds
via nearest-neighbor hopping through the faces of the PbI$_6$
octahedra, with a computed activation barrier of approximately
0.58~eV~\cite{Eames2015}.  Bromide vacancies in CsPbBr$_3$ have a
somewhat higher barrier of $\sim$0.74~eV, consistent with the greater
thermal stability of that compound.  Chloride vacancies, which are the
operative mobile species in MHyPbCl$_3$, have not been as extensively
studied computationally, but experimental activation energies inferred
from temperature-dependent dark current measurements suggest similar
or slightly higher barriers than for iodide.

\subsection{Consequences for Detector Performance}

Ion migration manifests in detector operation through three interrelated
phenomena.  First, under sustained bias, mobile ions drift to the
electrodes and accumulate, establishing an internal dipole that
\emph{screens} the applied field and reduces the effective field within
the crystal bulk.  This polarization effect depresses charge collection
efficiency over time and produces the characteristic dark current spike
followed by slow decay observed in the current transient of a freshly
biased perovskite crystal~\cite{He_NatPhoton_2021}.  Second, the
accumulation of ions at the electrode interfaces modifies the
contact barrier height, causing the rectification ratio and dark current
to evolve continuously during operation — a direct threat to the
baseline stability required in nuclear instrumentation.  Third, if
the bias polarity is reversed (as during AC-coupled readout or
impedance spectroscopy measurements), the accumulated ion profiles
transiently drive the internal field in the opposite direction to the
intended collection field, producing severe measurement artifacts.

These effects are not merely theoretical concerns: in a systematic study
of CsPbBr$_3$ single-crystal detectors, Pan et al.\ demonstrated that
the charge collection efficiency under a fixed field of
200~V\,cm$^{-1}$ dropped by approximately 30\% over the first
100~seconds of bias application before reaching a new quasi-equilibrium
dominated by ionic rather than electronic currents~\cite{Pan2019}.
The voltage-induced degradation was partially reversible upon removal
of bias, with recovery occurring on timescales of minutes to hours —
consistent with thermal relaxation of the accumulated ion
profiles~\cite{Zhao2019}.

\subsection{Mitigation Strategies}

Several strategies have been demonstrated to suppress or circumvent
ion-migration instability in perovskite detectors.  \emph{2D spacer layers}
(the Ruddlesden-Popper architecture) interpose wide-gap organic bilayers
at regular intervals through the crystal thickness, acting as
electrostatic barriers to halide vacancy drift while preserving charge
transport via tunneling or thermally activated hopping~\cite{Tsai2018}.
\emph{Surface passivation} with Lewis-base ligands such as
phenethylammonium iodide (PEAI), methylammonium chloride, or
thiophene-2-carboxylic acid coordinates unsaturated surface Pb$^{2+}$
sites that would otherwise serve as ion-migration nucleation points and
electronic trap states~\cite{Cao2018}.  \emph{All-inorganic compositions}
such as CsPbBr$_3$ exhibit significantly slower ion drift than their
hybrid organic-inorganic analogs because the rigid Cs$^+$ sublattice
provides fewer low-energy migration pathways.  Finally,
\emph{low-temperature operation} reduces the ionic conductivity
exponentially (Arrhenius behavior) and has been used to virtually
eliminate polarization artifacts in experimental detector
characterizations, though at the cost of the room-temperature operability
that is one of perovskite's primary advantages~\cite{He_NatPhoton_2021}.
Chapter~\ref{ch:2Dperovskite} investigates the combination of 2D
architecture with chemical surface passivation as a room-temperature
solution to the voltage-induced instability problem.

%% ------------------------------------------------------------
\section{Prior Art in Perovskite Radiation Detectors}
\label{sec:prior_art}
%% ------------------------------------------------------------

\subsection{Early Single-Crystal Gamma Detectors}

The first demonstration of halide perovskite single crystals as
room-temperature radiation detectors appeared in a pair of landmark
papers in 2016.  Wei et al.\ grew large MAPbI$_3$ single crystals by
inverse temperature crystallization and contacted them in an asymmetric
Au/perovskite/Cr geometry; at 662~keV they observed a resolved
$^{137}$Cs photopeak with an energy resolution of approximately
6.8\%~\cite{Wei2016}.  While this was inferior to the best CdZnTe
devices, the demonstration that a solution-processed hybrid perovskite
could rival polycrystalline CdZnTe — at a fraction of the material cost
and with far simpler growth conditions — generated immediate and sustained
interest.  Concurrently, Yakunin et al.\ reported sensitivity values
exceeding $10^3$~$\mu$C\,Gy$^{-1}$\,cm$^{-2}$ for MAPbI$_3$
thick-film X-ray detectors, establishing the perovskite
platform for medical and industrial X-ray imaging~\cite{Yakunin2015}.

\subsection{CsPbBr$_3$ and the Resolution Record}

The shift from hybrid to all-inorganic CsPbBr$_3$ was motivated by
the thermal and photo-instability of MAPbI$_3$, which loses iodide
at temperatures above $\sim$85~$^\circ$C and degrades rapidly in
ambient humidity.  Pan et al.\ demonstrated in 2019 that centimeter-scale
CsPbBr$_3$ single crystals grown by the Bridgman method exhibit a
resistivity exceeding $10^9$~$\Omega$\,cm, $\mu_e\tau_e$ products
above $10^{-2}$~cm$^2$\,V$^{-1}$, and a 662~keV energy resolution of
approximately 3.9\%~\cite{Pan2019}.  The subsequent integration of
CsPbBr$_3$ detectors with custom-designed ASICs — rather than
laboratory charge-sensitive amplifiers — addressed the electronic noise
limitation and enabled the record 1.4\% resolution at 662~keV reported
by He et al.\ in 2021~\cite{He_NatPhoton_2021}, subsequently approached
and extended by other groups employing refined crystal growth and improved
contact engineering~\cite{NatComm_CsPbBr3_ASIC}.

\subsection{X-Ray Detection and Imaging}

Polycrystalline thick films of MAPbI$_3$, MAPbBr$_3$, and
CsPbI$_2$Br have been investigated extensively as direct-conversion
X-ray detector arrays for applications in digital radiography and
computed tomography.  The high atomic number of the Pb sublattice
confers exceptional sensitivity at soft X-ray energies
(below $\sim$30~keV), where the photoelectric cross-section makes
lead-bearing compounds intrinsically more absorbing than the amorphous
selenium currently deployed in flat-panel detectors.  Wei et al.\ reported
a sensitivity of 4,500~$\mu$C\,Gy$^{-1}$\,cm$^{-2}$ for a 2-mm-thick
solution-processed MAPbI$_3$ film under an 8~keV X-ray beam~\cite{Wei2017},
and subsequent optimization of film deposition and contact passivation
has pushed sensitivities above 10$^4$~$\mu$C\,Gy$^{-1}$\,cm$^{-2}$
in several laboratories.  The principal challenges remaining for
commercial deployment are long-term operational stability (particularly
under the high photon fluences of medical imaging), the toxicity of
lead-bearing materials in a clinical environment, and the uniformity
of thick-film deposition over large detector areas.

\subsection{Neutron Detection}

Neutron detection by perovskite materials has been demonstrated via
two complementary approaches: neutron-capture via the incorporation
of $^{10}$B or $^6$LiI into the perovskite matrix, and direct fast-neutron
detection via hydrogen-recoil in hydrogen-rich perovskites.  The
composite approach was demonstrated in this group's own prior work:
MAPbI$_3$ crystals incorporating hexagonal boron nitride (hBN)
nanoflakes were shown to exhibit a measurable response to thermal
neutrons, with the neutron signal distinguished from the ambient gamma
background by pulse-shape analysis~\cite{Panaccione_NIM_2024}.  The
MHyPbCl$_3$ results reported in Chapter~\ref{ch:MHyPbCl3} represent
the first systematic investigation of a perovskite designed from the
outset for high hydrogen content, aiming to maximize the fast-neutron
elastic-scatter cross-section in a single-phase semiconducting medium.
The availability of a structurally well-characterized perovskite with
both high H density and a bandgap suitable for room-temperature
electronic detection opens a path toward compact direct-conversion
fast-neutron detectors that do not require external moderators or
$^3$He-filled proportional tubes.

\subsection{Mixed-Cation and 2D Detector Architectures}

The instability of single-phase 3D perovskites under sustained bias
has driven interest in compositionally engineered architectures designed
to suppress ion migration while retaining the favorable electronic
properties of the perovskite host.  Two-dimensional Ruddlesden-Popper
compounds were first applied to radiation detection by Xu et al., who
demonstrated that PEA$_2$MA$_3$Pb$_4$I$_{13}$ ($n = 4$) single crystals
exhibited suppressed polarization current compared to 3D MAPbI$_3$
under identical bias conditions, attributed to the PEA$^+$ barriers
impeding halide vacancy drift~\cite{Xu2020}.  The compositions studied
in Chapter~\ref{ch:2Dperovskite} extend this concept to the $n = 3$
compound PEA$_2$MA$_2$Pb$_3$I$_{10}$, where the additional organic
layer density provides even stronger suppression of ion migration, and
investigates the interaction between the 2D architecture and chemical
surface passivation as a dual strategy for achieving stable, low-dark-current
detector operation.

%% ------------------------------------------------------------
\section{Summary}
%% ------------------------------------------------------------

This chapter has established the conceptual framework necessary to
interpret the experimental results reported in Chapters~\ref{ch:CsFAPbI}
through~\ref{ch:MHyPbCl3}.  The ABX$_3$ halide perovskite framework
accommodates a diverse family of compositions spanning 3D, quasi-2D, and
layered architectures, with A-site chemistry tunable from inorganic Cs$^+$
through hybrid organic-inorganic MA$^+$ and FA$^+$ to the
hydrogen-rich methylhydrazinium cation of MHyPbCl$_3$.  The physics of
semiconductor radiation detectors connects material properties —
bandgap, effective atomic number, $\mu\tau$ product — to measurable
figures of merit such as energy resolution, sensitivity, and neutron
detection efficiency via the Hecht equation and Klein's pair-creation
rule.  The Schottky diode contact architecture suppresses dark current
to levels compatible with radiation counting without sacrificing charge
collection, provided that ion migration is controlled by compositional
engineering or surface passivation.  The problem of voltage-induced
instability driven by halide vacancy migration is identified as the
central unsolved challenge for room-temperature perovskite detector
operation, motivating the 2D and passivation strategies pursued in
the experimental chapters that follow.